% Fuente: Auxiliar 2, P1 (pauta).
\textbf{Variables de decisión}

Se busca minimizar los costos de habilitar enlaces y costos de transporte para abastecer la demanda. Por lo tanto, las variables de decisiones serán:

\begin{equation}
y_{i,j} =
\begin{cases}
1 & \text{si se habilita el enlace de } i \text{ a } j, \\
0 & \text{caso contrario}
\end{cases}
\end{equation}

\begin{equation}
x_{i,j} \;:\; \text{Flujo de } i \text{ a } j
\end{equation}

El dominio de la variable $x_{i,j}$ serán los reales positivos debido a que es un grafo dirigido
(el flujo puede tomar una sola dirección). Notar que en el caso de que se requiera un flujo
bidireccional de $i$ a $j$ se requiere habilitar los enlaces $x_{i,j}$ y $x_{j,i}$.\\

\textbf{Parámetros}

\begin{itemize}
    \item $N$: Nodos.
    \item $c_{i,j}$: Costos por unidad de flujo.
    \item $d_{i,j}$: Costos por construcción.
    \item $U_{i,j}$: Capacidad máxima del enlace.
    \item $b_n$: Demanda/Suministro del nodo $n$, asumiremos que es positiva si corresponde a demanda y negativa en el caso de ser suministro.
\end{itemize}

\textbf{Función objetivo}

\begin{equation}
\min_{x_{i,j},\,y_{i,j}} \sum_{(i,j)\in A} c_{i,j}\,x_{i,j} \;+\; \sum_{(i,j)\in A} d_{i,j}\,y_{i,j}
\end{equation}

donde el primer término corresponde a los costos por transporte de flujo, mientras que el segundo corresponde a los costos por construcción de enlace.\\

\textbf{Restricciones}

\begin{itemize}
    \item[a)] Balance: Los flujos entrantes a cualquier nodo $n \in N$ deben ser iguales a los flujos salientes más la cantidad demandada o suministrada.

    \begin{equation}
    \sum_{j\,|\,(n,j)\in A} x_{n,j} + b_n = \sum_{i\,|\,(i,n)\in A} x_{i,n}
    \qquad \forall n \in N
    \end{equation}

    \item[b)] Límites de flujo: Se deberá respetar la capacidad máxima de flujo por el enlace. Adicionalmente, flujos por enlaces podrán ser distintos de cero solo en caso de decidir ser habilitados ($y_{i,j}=1$).
    
    \begin{equation}
    x_{i,j} \leq U_{i,j}\,y_{i,j} \qquad \forall (i,j) \in A
    \end{equation}

    \item[c)] Dominio de las variables:

    \begin{equation}
        x_{i,j} \geq 0 \qquad \forall (i,j) \in A
    \end{equation}

    \begin{equation}
        y_{i,j} \in \{0,1\} \qquad \forall (i,j) \in A
    \end{equation}
    
\end{itemize}
